\documentclass[11pt]{article}
\usepackage[margin=1in]{geometry}
\usepackage{amsmath,amssymb,microtype}
\usepackage[colorlinks=true,urlcolor=blue,citecolor=blue]{hyperref}
\title{Applebaum's stable trace-asymptotic conjecture was answered by arXiv:1308.4944}
\author{Literature-status note for \texttt{fa\_banach\_001}}
\date{26 August 2026}

\begin{document}
\maketitle

\noindent\textbf{Status.}
\texttt{literature\_already\_answered (full conjecture)}.  No new
mathematical result is claimed by this packet.

\section*{Original conjecture}

Let $G$ be a compact semisimple Lie group of dimension $d$.  For the central
$\alpha$-stable semigroup constructed from the Casimir operator, Applebaum
uses the exponent $\eta(u)=|u|^\alpha$ and denotes its continuous density by
$k_t$.  In Section 6, arXiv PDF page 21 (journal page 2494), he conjectures
that, for $0<\alpha\leq 2$,
\begin{equation}\label{eq:applebaum}
  k_t(e)\sim C t^{-d/\alpha}\qquad(t\downarrow0),
\end{equation}
with $C>0$; the Cauchy conjecture is the case $\alpha=1$.  The same section
has already established $k_t(e)=\operatorname{Tr}(T_t)$.

\section*{Explicit later answer}

Banuelos and Baudoin state in the abstract and first paragraph of
arXiv:1308.4944, PDF page 1, that they address Applebaum's conjecture and that
their purpose is to show it is indeed true \cite{BB2013}.  They cite
Applebaum's paper as reference [1] and identify the conjecture at its journal
pages 2493--2494.

Their pages 2--3 take
\[
  \psi(\lambda)=\lambda^{\alpha/2},\qquad 0<\alpha<2,
\]
so that the subordinate generator is $(-\Delta)^{\alpha/2}$.  If $M$ is a
compact Riemannian $n$-manifold, Weyl's law and a Tauberian argument give
\begin{equation}\label{eq:bb-stable}
 \operatorname{Tr}(Q_t)
 \sim
 \frac{\operatorname{Vol}(M)\,\Gamma(n/\alpha+1)}
      {\Gamma(n/2+1)(4\pi)^{n/2}}\,t^{-n/\alpha}.
\end{equation}
The unnumbered theorem and equation (9) on PDF page 3 prove a more general
regular-variation version.  Immediately after (9), they observe that on a
compact Lie group with normalized Haar measure the trace asymptotic gives the
on-diagonal kernel asymptotic.  Applying this to the bi-invariant Laplacian on
$G$ yields \eqref{eq:applebaum}.  This is exactly Applebaum's construction,
because $\eta(\sqrt{\kappa})=\kappa^{\alpha/2}$ on every Casimir eigenspace.

The answering note treats $0<\alpha<2$.  The endpoint $\alpha=2$ is the
classical heat-kernel case, which Applebaum explicitly records as already
known in the same section.  Consequently the combined literature answers the
full stated range $0<\alpha\leq2$, and even proves the trace result on every
closed Riemannian manifold.

\section*{Search and scope}

The four run indexes, exact-title and DOI hits in the local parsed arXiv
corpus, exact-phrase web searches, and the 16 OpenAlex-indexed citing works
were checked through 26 August 2026.  arXiv:1308.4944 is a direct, explicit
answer rather than an agent-identified implication.  Later papers by
Fahrenwaldt and by Applebaum--Le Ngan corroborate its use.  No portion of
Applebaum's stable-asymptotic conjecture is claimed open here.

\begin{thebibliography}{9}
\bibitem{Applebaum2011}
D.~Applebaum,
\emph{Infinitely divisible central probability measures on compact Lie
groups---regularity, semigroups and transition kernels},
Ann. Probab. \textbf{39} (2011), 2474--2496;
arXiv:1006.4711,
\url{https://arxiv.org/abs/1006.4711}.

\bibitem{BB2013}
R.~Banuelos and F.~Baudoin,
\emph{Trace asymptotics for subordinate semigroups},
arXiv:1308.4944 (2013),
\url{https://arxiv.org/abs/1308.4944}.
\end{thebibliography}

\end{document}
